\section{怎样辨析反义词}

读报时，我们会碰到这样一类的句子：

“在战争中，我们首先要调查清楚敌我双方的虚实，强弱以及士气的高低或低落，然后才能决定进退、攻守以及全面包围或局部分割。”

这句话中的“虚”与“实”、“强”与“弱”、“高昂”与“低落”、“进”与“退”、“攻”与“守”，都是意义相反的词，它们之间的意义互相排斥；“全面”与“局部”、“包围”与“分割”、“敌”与“我”，都是意义相对的词，它们之间的意义互相对立，并有别的相对概念。语言学家把语言中这一类意义相反或者意义相对的词叫做反义词。

怎样辨析反义词呢？以下几点请加注意：

1. 认真辨析这些词的意义是否相反或者相对。除了上述例句中的一些反义词之外，象“长 --- 短”、“大 --- 小”、\\
“好 --- 坏”、“出 --- 入”、“买 --- 卖”、“公 --- 私”、“古 --- 今”、“正确 --- 错误”、“伟大 --- 渺小”、“勇敢 --- 怯懦”、“完整 --- 残缺”、“团结 --- 分裂”、“胜利 --- 失败”、\\
“建设 --- 破坏”、“拥护 --- 反对”、“虚心 --- 骄傲”、“真理 --- 谬误”、“现象 --- 本质”、“经验 --- 教训”、“白天 --- 黑夜”、“唯物论 --- 唯心论”，等等，都是意义相对或者相对的反义词，都反映了客观事物的矛盾对立的现象。

有些词就意义来说，并没有明显的相反相对的关系，但是在一定的上下文里却表示了相反或相对的意义，这样用时，也可以把它们看成是反义词（广义）。例如：

“我们的痈疽，是它们的宝贝，那么，它们的敌人，当然是我们的朋友了。”（鲁迅：《南腔北调集·我们不再受骗了》）

“敌人”和“朋友”是意义对立的反义词，而“痈疽”的意思是指毒疮，本来并不能与“宝贝”构成反义词，这儿却把二者进行对举，可以看成是广义的反义词。

另外，由于语言习惯，人们经常把某些意义不是相反或者相对的词，如“春”与“秋”、“黑”和“白”，等并举对比，也应把它们看成是反义词。

除了语言习惯或者上下文中显示出相反、相对的意义的词之外，我们把意义不是相反或相对的词，都不看成是反义词。在语言运用中，用否定词语也能表现事物之间矛盾对立的关系，例如“好”与“不好”，“坚定”和“不坚定”之类，也不能说它们是反义词，因为“不好”、“不坚定”都不是一个词，而是一个词组了。

按照这种辨析方法，我们可以给常见词（主要是形容词、动词和名词）找出它的反义词。例如：“吝啬 --- 慷慨”，“崎岖 --- 平坦”，“前进 --- 后退”，“糟粕 --- 精华”，\\
“形象 --- 抽象”，等等。也可以完成一些前后词语结构相似、意义相反或相对的成语。例如上行（下效），朝秦（暮楚），隐恶（扬善），无独（有偶），博古（通今），口蜜（腹剑），舍近（求远），阳奉（阴违），争先（恐后），死去（活来），弃暗（投明），惩前（毖后），醉生（梦死），朝发（夕至），头重（脚轻），承前（启后），顺之者昌（逆之者亡），兼听则明（偏信则暗），等等。

有些词的确很难找到和它相反、相对的反义词，象“彻底”、“打算”等，不过这毕竟是个别的现象，并不影响我们对于反义词的辨析。

2. 认真辨析词性是否相同。一般说来，反义词的词性是相同的。例如开头所述例句中的“虚 --- 实”、“强 --- 弱”、“高昂 --- 低落”、“全面 --- 局部”，在这儿都是形容词；“进 --- 退”、“攻 --- 守”、“包围 --- 分割”，都是动词；“敌（军） --- 我（军）”则是名词或代名词。可见，反义词以形容词、动词以及名词为常见。

如果词性不同，一般是不能构成反义词的。例如：

“国土的完整或分裂，关系到国家的统一或残缺。”

我们只要细心地辨析一下，就会发现这个句子中的反义词用得不对。你看，“完整”是形容词，而“分裂”是动词，它们怎么会是反义词呢？“统一”（动词）与“残缺”（形容词）之间同样不能构成反义词。我们如果把句中的“分裂”与“残缺”两个词对调一下，使“完整”与“残缺”、“统一”与“分裂”构成反义词，这个句子就通顺了。

根据这种辨析方法，我们就可以避免或者改正由于不明确反义词要求词性一致的常识而出现的病句。例如：

（1）民间故事大都是歌颂劳动人民的善良和机智，揭露敌人的暴行和愚蠢的。

（“暴行”是名词，不能与形容词“善良”构成反义词，应改为“丑恶”或“残暴”。）

（2）对于违犯党纪国法的，不管是什么人，都要依法处理，以伸张正气，消除邪恶。

（“邪恶”是形容词，不能与名词“正气”构成反义词，应改为“邪气”。）

3. 认真辨析反义词的对应关系。由于客观事物的矛盾对立现象是错综复杂的，如果是一个多义词，就可能有几个不同的反义词。例如：

（1）虚心使人进步，骄傲使人落后。

（2）作为一个新中国的青年，我感到非常骄傲。

第一句中的“骄傲”是“自满”的意思，反义词就是“虚心”；第二句中的“骄傲”是“自豪”的意思，反义词只能是“自卑”了。

再如：“公”的反义词可以是“私”，也可以是“母”；“实”的反义词可以是“虚”，也可以是“空”；“来”的反义词可以是“去”，也可以是“往”；“老”的反义词可以是“幼”、“少”，也可以是“新”，还可以是“嫩”；“深”的反义词可以是“浅”、“淡”或者“薄”；“开”的反义词可以是“合”、“关”或者“停”；“去”的反义词可以是“来”、“留”或者“就”；“分裂”的反义词可以是“团结”或者“统一”；“进步”的反义词可以是“倒退”、“落后”或者“反动”；“反对”的反义词可以是“拥护”、“支持”或者“赞成”，等等。这就是说，这一类的词都是多义词。而在不同的句子中，由于一定的上下文决定了某个多义词的特定意义，因而它的反义词也就被决定下来了。例如“进步”这个多义词，在句子中是“前进”的意思时，它的反义词是“倒退”，是“向前发展”的意思时，它的反义词是“落后”；是“革命”的意思时，它的反义词则只能是“反动”了。

我们在辨析反义词时，必须联系一定的上下文，弄清这个词在句子中的含义，这样做了之后，与它对应的反义词就可以确定下来了。

4. 认真辨析反义词的修辞作用。由于反义词能鲜明地揭示出矛盾事物的两个对立面，更清楚地暴露矛盾事物的对立性，所以在句子中反义词经常被用作修辞上对比、映衬的手段。例如：“旧社会把人逼成鬼，新社会将鬼变成人。”“人无远虑，必有近忧。”“有这个借鉴和没有这个借鉴是不同的，这里有文野之分，粗细之分，高低之分，快慢之分。”这样运用反义词，能使文章具有更加鲜明的色彩和更加强烈的说服力。
\newpage

需要注意的是，有些文章常利用反义词意义的矛盾对立，来构成“反语”。因此，我们对这种作为修辞手法之一的“反语”，只能从反面去理解其意义，切不可正面看。例如用“小鬼头”来称呼所喜爱的孩子，用“你真坏”来表示对爱人的由衷喜爱的感情，意味更加深厚。相反，用含有褒义的反义词来代替原来的贬义词，则构成一种更深刻的否定和讥讽的感情意味。例如在毛泽东同志的著作中，毛泽东同志常用“亲爱的国民党先生们”、“可爱的蒋总统”等词语，表达了对国民党反动派及其头子的极度憎恶和嘲笑的强烈感情。

5. 注意掌握反义词的某些固定用法。汉语中可以利用

反义词来构成合成词，成语、谚语，而其中所用的反义词都应遵循约定俗成的规则，不得随意更改。例如利用反义词构成的合成词“是非”（不能改为“对非”）、“朝夕”（不能改为“朝晚”）、“毁誉”（不能改为“诽谤”）等。

对于成语中的反义词，尤其不能随意更换，因为成语是固定词组。例如：“化险为夷”（“夷”不能写成“安”），“无独有偶”（“偶”不能写成“双”），色厉内荏（“荏”不能写成“弱”），“通宵达旦”（“宵”不能写成“夜”），言简意赅（“赅”不能写成“全”），世态炎凉（“炎”不能写成“暖”），“一暴十寒”（“暴”不能写成“晒”），“志大才疏”（“疏”不能写成“浅”），“欲擒故纵”（“纵”不能写成“放”），一劳永逸（“逸”不能写成“安”），“生杀予夺”（“予”不能写成“与”），新陈代谢（“陈”不能写成“旧”），宁缺勿滥（“缺”不能写成“少”），等等，都是固定搭配，如果拿另外同义或者近义的字去替换其中的任何一个反义词，那是不行的。
\newpage

至于谚语，则是人民群众口头上流传的通俗、简炼、含义深刻的现成话，读起来顺口，听起来悦耳，易懂易记，口耳相传，也具有约定俗成的特点，不能随意更改。例如：“富人一席酒，穷汉半年粮”、“地是黄金板，人勤地不懒”、“水往低处流，人往高处走”、“三个臭皮匠，凑成诸葛亮”、“只许州官放火，不许百姓点灯”、“真的难灭，假的易除。不经冬寒，那知春暖”，等等，它们中的反义词，以鲜比的对比、映衬的方式，把抽象而富有哲学意味的道理表达出来，引起人们的联想，使人得到深刻的印象。我们在运用这些谚语时，当然也不能随意地加以改变。
\newpage